\documentclass[a4paper,12pt]{article}
\usepackage{amsmath, amssymb}
\usepackage{xcolor}
\usepackage{tcolorbox}
\usepackage{multicol}
\usepackage{geometry}
\hyphenpenalty=100000

\geometry{left=1.5cm, right=1.5cm, top=2cm, bottom=2cm}

\tcbuselibrary{theorems}

\definecolor{PurpleEx}{HTML}{5d478b}
\definecolor{GrayEx}{HTML}{f2f2f2}
\definecolor{GrayBG}{HTML}{d9d9d9}
\definecolor{BlackSav}{HTML}{000000}

\title{
\centering
\tcbset{colframe=BlackSav,colback=GrayBG}
\tcbox[tikznode]{
Chapitre 3 \\ Devoir maison 1
}
}
\usepackage{tikz}
\author{}
\date{}

\newtcbtheorem
  []% init options
  {exercice}% name
  {Exercice}% title
  {%
    colback=GrayEx,
    colframe=PurpleEx,
    fonttitle=\bfseries,
  }% options
  {ex}% prefix


\begin{document}

\maketitle

% Exercice 18
\begin{exercice}[title={Exercice 18 - {[}Représenter, Calculer, Raisonner{]}}]{}{}
\begin{enumerate}
    \item On dit qu'un triangle équilatéral $ABC$ est direct si et seulement si $(\overrightarrow{AB}; \overrightarrow{AC}) = \frac{\pi}{3}$.
    Soient \( A \), \( B \) et \( C \) trois points non alignés d’affixe \( a \), \( b \) et \( c \). On note \( j = e^{\frac{2i\pi}{3}} \).
    \begin{enumerate}
        \item \begin{enumerate}
            \item Vérifier que $1, j$ et $j^2$ sont les solutions de l'équation $z^3 = 1$.
            \item Calculer $(1-j)(1 + j + j^2)$. En déduire que $1 + j + j^2 = 0$.
            \item Vérifier que
            $$e^{i\frac{\pi}{3}} + j^2 = 0$$
        \end{enumerate}
        \item \begin{enumerate}
            \item Démontrer que le triangle $ABC$ est équilatéral direct si et seulement si :
            $$\frac{c - a}{b - a} = e^{i\frac{\pi}{3}}$$
            \item Montrer que le triangle \( ABC \) est équilatéral direct si, et seulement si,
            \[
            a + bj + cj^2 = 0.
            \]
        \end{enumerate}
    \end{enumerate}
    \item Soit le centre de gravité d'un ensemble de $n$ points d'affixes $z_1, \dots z_n$ le point d'affixe
    $$\dfrac{z_1 + \dots + z_n}{n}$$ 
    \begin{enumerate}
        \item On ne suppose pas nécessairement que \( ABC \) est équilatéral. On construit à partir de \( ABC \) les trois triangles équilatéraux de base \([AB]\), \([AC]\) et \([BC]\) construits à l’extérieur de \( ABC \). Montrer que les centres de gravité de ces trois triangles forment un triangle équilatéral.
    \end{enumerate}
\end{enumerate}
\end{exercice}

\end{document}